\documentclass[aspectratio=169]{beamer}
\usepackage{sudsley}

% A UML stick figure for use-case diagrams: \actor{(x,y)}{name}{label}.
\newcommand{\actor}[3]{%
  \begin{scope}[shift={#1}]
    \draw (0,0.42) circle (0.12);
    \draw (0,0.30) -- (0,0) (-0.2,0.2) -- (0.2,0.2) (0,0) -- (-0.15,-0.25)
          (0,0) -- (0.15,-0.25);
    \node[font=\small, below] at (0,-0.28) {#3};
    \coordinate (#2) at (0,0.15);
  \end{scope}}

\title{System Design}
\subtitle{Introduction}

\begin{document}

\maketitle

\begin{frame}{What is system design?}
  Also called architecture design. It covers:
  \begin{columns}[T]
    \begin{column}{0.48\textwidth}
      \small
      \begin{itemize}
        \item Functional and non-functional requirements
        \item Capacity estimation
        \item Relational and NoSQL databases
        \item Vertical scaling, horizontal scaling, sharding
        \item Load balancing
        \item Primary--secondary replication
        \item Caches and CDNs
        \item Stateless and stateful servers
      \end{itemize}
    \end{column}
    \begin{column}{0.48\textwidth}
      \small
      \begin{itemize}
        \item Datacenter georouting
        \item Message queues, publish--subscribe
        \item Performance metrics, monitoring, logging
        \item Build, test, configure, and deploy automation
        \item Finding single points of failure
        \item API rate limiting
        \item Service level agreements
      \end{itemize}
    \end{column}
  \end{columns}
\end{frame}

\begin{frame}{Resources for system design interviews}
  \begin{enumerate}
    \item Sukumar Ghosh, \emph{Distributed Systems: An Algorithmic Approach}
    \item The System Design Primer:
      \url{https://github.com/donnemartin/system-design-primer}
    \item Roberto Vitillo, \emph{Understanding Distributed Systems}, 2nd ed.
  \end{enumerate}
\end{frame}

\begin{frame}{How we'll learn}
  \centering
  \begin{tabularx}{\textwidth}{@{}>{\bfseries}l X@{}}
    \toprule
    Lecture & Topics \\
    \midrule
    1. Introduction & The scenario, product ideas, the MVP, high-level design \\
    2. High-level and component design & The model design, the thirteen component questions \\
    3. Component implementation & Approaching a new codebase \\
    4. Code review and analysis & Reading code, the cost of code \\
    \bottomrule
  \end{tabularx}
\end{frame}

\section{Welcome to Sudsley!}

\begin{frame}{Welcome to Sudsley!}
  \Large Today, you're joining the founding engineering team of a disruptive new startup.
\end{frame}

\begin{frame}{The problems we're solving}
  \begin{enumerate}
    \item<1-> \textbf{Wasted time.} Customers travel to and wait at the laundromat, unsure
      when their machine will be done.
    \item<2-> \textbf{Inconvenience.} Machines are occupied, forcing customers to wait or
      come back later.
    \item<3-> \textbf{Forgotten clothes.} Clothes are often left in machines long after the
      cycle finishes.
  \end{enumerate}
\end{frame}

\begin{frame}{Our secret weapon: the smart laundry machine}
  At the heart of Sudsley is our proprietary, next-generation washer and dryer. Some very
  smart engineers spent years developing it.

  \bigskip
  Every machine is internet-enabled, with an integrated API for complete remote control and
  monitoring:
  \begin{itemize}
    \item \texttt{startCycle()}
    \item \texttt{stopCycle()}
    \item \texttt{getStatus()}: available, running, finished, time remaining, needs
      maintenance
  \end{itemize}
\end{frame}

\begin{frame}{Your job}
  \begin{enumerate}
    \item Explore product options and suggest ideas.
    \item Design the product.
    \item Implement the core components.
  \end{enumerate}
\end{frame}

\section{Product ideas}

\begin{frame}{Product ideas}
  \Large What software should we build, and why?

  \bigskip
  \normalsize\alert{Discuss in groups.}
\end{frame}

\begin{frame}{Great ideas! But they don't matter that much}
  The CEO and product managers (with your completely relevant input) have settled on an idea:
  \alert{UberEats, Postmates, or DoorDash, but for laundry.}

  \medskip
  How they want it to work:
  \begin{enumerate}
    \small
    \item A customer puts their laundry out on their front porch.
    \item They open an app to request a launderer.
    \item The launderer picks up the clothes and takes them to a Sudsley location.
    \item The launderer puts the laundry in a machine.
    \item The app keeps the customer up to date on the status of their clothes.
    \item The launderer delivers the clothes back once they're finished.
    \item The launderer gets a portion of the job revenue.
  \end{enumerate}
\end{frame}

\section{Minimum viable product}

\begin{frame}{Define the MVP requirements}
  The \alert{minimum viable product} is an early version of a product with just enough
  features to be functional.

  \medskip
  Guidelines for an MVP:
  \begin{itemize}
    \item Identify the must-have features.
    \item Avoid unnecessary work.
    \item Expect feedback from early users. This is not the forever-system.
  \end{itemize}

  \medskip
  \alert{In groups:} define a rough MVP for the alpha version of Sudsley.
\end{frame}

\begin{frame}{Must-haves}
  \begin{columns}[T]
    \begin{column}{0.5\textwidth}
      \small
      \begin{itemize}
        \item Assume about 10k initial customers
        \item Customers can request a laundry job
        \item Launderers can accept a laundry job
        \item Launderers know which machine to put the clothes in
        \item Customers know where their clothes are geographically
      \end{itemize}
    \end{column}
    \begin{column}{0.5\textwidth}
      \small
      \begin{itemize}
        \item Customers know the status of their laundry job
        \item Customers pay through the app
        \item Launderers get paid through the app
        \item Customers can rate their laundry job
        \item Basic statistics and telemetry
      \end{itemize}
    \end{column}
  \end{columns}
\end{frame}

\section{High-level design}

\begin{frame}{Define the high-level design}
  High-level (or macro-level) design is the first formal step of system design. It should be
  easy to follow.
  \begin{itemize}
    \item Identifies the necessary components (microservices, external services, clients)
      without detailing their implementation
    \item Aligns everyone on the project's goals and helps distribute responsibilities
  \end{itemize}
  May also include:
  \begin{itemize}
    \item UI designs and data flow diagrams
    \item The technology stack
    \item The deployment procedure
  \end{itemize}
\end{frame}

\begin{frame}{Design tools}
  Many high-level designs are drawn as UML diagrams, though this is not strictly required.
  \begin{itemize}
    \item draw.io: \url{https://app.diagrams.net/}
    \item Lucidchart: \url{https://www.lucidchart.com/}
    \item Mermaid: \url{https://mermaid.js.org/}
  \end{itemize}
\end{frame}

\begin{frame}{Example: a use-case diagram}
  \centering
  \begin{tikzpicture}[
      uc/.style={ellipse, draw, minimum width=28mm, minimum height=9mm, font=\small},
      every edge/.style={draw, thin}]
    \node[draw, rounded corners, minimum width=48mm, minimum height=50mm,
      label={[font=\small]above:Online bookstore}] (sys) {};
    \node[uc] (browse) at ([yshift=16mm]sys.center) {Browse books};
    \node[uc] (buy) at (sys.center) {Buy a book};
    \node[uc] (login) at ([yshift=-16mm]sys.center) {Log in};
    \actor{(-5,0)}{cust}{Customer}
    \actor{(5,0)}{pay}{Payment provider}
    \actor{(5,-1.6)}{idp}{Identity provider}
    \path (cust) edge (browse.west) (cust) edge (buy.west) (cust) edge (login.west)
          (buy.east) edge (pay) (login.east) edge (idp);
    \draw[-{Stealth}, dashed] (buy) -- node[right, font=\scriptsize]{include} (login);
  \end{tikzpicture}
\end{frame}

\begin{frame}{Homework 1: high-level design (collaborative)}
  \begin{enumerate}
    \item Based on the product and MVP requirements from class, write a brief description of
      how you think the system should look and operate. No more than half a page.
    \item Create a high-level design for the Sudsley system.
      \begin{itemize}
        \item A business or system use-case diagram is a good start.
        \item You don't need to follow UML strictly. Make it readable and intuitive.
      \end{itemize}
  \end{enumerate}
  \alert{Be flexible! No design is correct or incorrect.}
  \handout{assignments/hw1-high-level-design/HOMEWORK.md}
\end{frame}

\end{document}
